\documentclass[tikz,border=4pt]{standalone}
\usepackage{pe-tikz}
\begin{document}
\begin{tikzpicture}[v/.style={pe/node, minimum size=9mm, font=\small}, ar/.style={pe/flow}]
  \foreach \i/\name/\title/\desc in {0/mcar/MCAR/{a hiányzás független\\az adatoktól},
                                     1/mar/MAR/{a hiányzás a megfigyelt\\$X$-től függ},
                                     2/mnar/MNAR/{a hiányzás magától\\a hiányzó $Y$-tól függ}}{
    \begin{scope}[xshift=\i*5.2cm]
      \node[v] (X\i) at (0,1.4) {$X$};
      \node[v] (Y\i) at (2.4,1.4) {$Y$};
      \node[v, fill=pelight] (R\i) at (1.2,0) {$R$};
      \draw[pe/thin, draw=muted] (X\i) -- (Y\i);
      \node[font=\small\bfseries, text=peblue] at (1.2,2.4) {\title};
      \node[pe/note, text width=40mm] at (1.2,-1.1) {\desc};
    \end{scope}}
  \draw[ar] (X1) -- (R1);
  \draw[ar] (Y2) -- (R2);
  \draw[ar, draw=muted, dashed] (X2) -- (R2);
  \node[pe/note, text width=150mm, align=center] at (6.4,-2.0) {$X$: mindig megfigyelt változó,\enspace $Y$: hiányos változó,\enspace $R$: a hiányzás jelzője ($R=1$, ha $Y$ hiányzik)};
\end{tikzpicture}
\end{document}
